\chapter{Exponential and Logarithm}\label{ch:exponential-logarithm}

These functions are useful tests of the distinction between defining a
computation and proving equivalence with another one.  We construct the
exponential by factorial prefixes, the logarithm by rectangles, and then
prove that they are inverse computations.

\section{Factorial prefixes}
For a rational bound $R\ge|x|$, set
\[
 E_N(x)=\sum_{n=0}^N\frac{x^n}{n!}.
\]
If $N+2\ge2R$, the absolute terms after the first omitted one have
successive ratio at most $1/2$.  Thus
\[
 |E-E_N|\le2\frac{R^{N+1}}{(N+1)!}.
\]
Finite-prefix stabilization constructs $E(x)$ on every bounded interval.
The differentiated series is identical, so the previous chapter proves
$DE=E$, with a uniform increment error bound on each such interval.

Finite binomial algebra in the double series gives
\[
 E(x+y)=E(x)E(y),\qquad E(0)=1,\qquad E(x)E(-x)=1.
\]
To justify the rearrangement, first restrict to finitely many terms;
the omitted terms are bounded by the product of the factorial majorants.
For $x\ge0$, the coefficients give $E(x)\ge1+x$.  The reciprocal
identity then proves positivity for negative $x$ as well.

The midpoint difference has the form
\[
 E(x+h)+E(x-h)-2E(x)
   =E(x)\bigl(E(h)+E(-h)-2\bigr)\ge0,
\]
since the final factor is a sum of nonnegative even powers.  Hence $E$ is
convex.  The series remainder proves that its convex derivative computation
terminates and has value $E$.

\section{Compound interest is another computation}
For $N\ge1$, the binomial coefficient of $x^k/k!$ in $(1+x/N)^N$ is
\[
 q_{N,k}=\prod_{j=0}^{k-1}(1-j/N)\quad(k\le N),
\]
and it is zero for $k>N$.  For $k\le N$,
$0\le1-q_{N,k}\le k(k-1)/(2N)$, by expanding the finite product or
inductively using $1-uv\le(1-u)+(1-v)$ for $u,v\in[0,1]$.
For $k>N$ the same upper bound for $1-q_{N,k}=1$ still holds.
Consequently, if $|x|\le R$ and $B_R$ is a rational upper bound for $E(R)$,
\[
 \left|(1+x/N)^N-E(x)\right|
 \le\frac1{2N}\sum_{k\ge2}\frac{k(k-1)R^k}{k!}
 \le\frac{R^2B_R}{2N}.
\]
Every infinite sum in this estimate has already been given a factorial
tail.  This proves the compound-interest representation by a finite error
budget.  For its usual positive-base interpretation also choose $N>R$.

\section{The logarithm by accumulation}
For rational $x>0$, define
\[
 L(x)=\int_1^x\frac{dt}{t}.
\]
On every positive rational interval $[a,b]$, the integrand is decreasing,
bounded, and has the algebraic continuity estimate
$|1/x-1/y|\le |x-y|/a^2$.  The integral is therefore constructed
independently of $E$.  The accumulation estimate proves
\[
 DL(x)=1/x.
\]
Ordered integral averages make $L$ concave.  Its one-sided secants converge
by the same continuity estimate, so this is also its computed derivative.

For $0<x<y$ in $[a,b]$,
\[
 \frac{y-x}{b}\le L(y)-L(x)\le\frac{y-x}{a}.
\]
This supplies continuity and strict separation.  Rational scaling of
partitions gives, for positive rational $x,y$,
\[
 L(xy)=L(x)+L(y).
\]
The same identity extends to computed positive inputs by the
computed-input construction and the displayed bounds.  In particular
$L(2^n)=nL(2)$, with $L(2)\ge1/2$, gives explicit arbitrarily large
values, and reciprocation gives arbitrarily negative ones.

\section{Proving the inverse relation}
On a fixed bounded input interval choose rational bounds $B_E,K_E$ such
that $|E|\le B_E$ and
\[
 |E(x+h)-E(x)-E(x)h|\le K_Eh^2/2.
\]
The factorial tails and the series remainder supply them.  Also enclose
the image in $[c,d]$ with rational $c>0$, using $E(-x)=1/E(x)$.
Accumulation and the Lipschitz bound $c^{-2}$ for $1/t$ give
\[
 |L(z+k)-L(z)-k/z|\le k^2/(2c^2)
\]
for a segment in that image interval.  The derivative bound gives
$|E(x+h)-E(x)|\le B_E|h|$ (enlarge $B_E$ on the interval if necessary).
Substitution into these two estimates proves the explicit composite error
\[
 |L(E(x+h))-L(E(x))-h|
 \le\left(\frac{K_E}{2c}+\frac{B_E^2}{2c^2}\right)h^2.
\]

Adding these estimates on equal rational partitions shows that
$L(E(x))-x$ has endpoint difference at most
$\varepsilon$ times the interval length for every $\varepsilon>0$.
At zero it is zero.  Hence
\[
 L(E(x))=x.
\]
This is a finite telescoping proof, not the invalid implication
``pointwise rational derivative zero implies constant.''

The identities extend to computed inputs on bounded intervals by
continuity.  In particular, for a positive computed $y$, apply the identity
to $x=L(y)$.  Both $E(L(y))$ and $y$ have logarithm $L(y)$, and the
separation bound for $L$ proves their equivalence:
\[
 E(L(y))=y.
\]
Thus $E=\exp$ and $L=\log$ name inverse computations.

\section{Representative consequences}
Define $x^\alpha=E(\alpha L(x))$ on a positively separated interval,
for a supplied computed constant $\alpha$.  The same finite increment
composition gives
\[
 D(x^\alpha)=\alpha x^{\alpha-1}.
\]
This statement includes the domain and bounded-interval estimates; it is
not an assertion that arbitrary rational-point differentiable functions
can be freely composed and integrated.

A particularly useful primitive is $xL(x)-x$.  The product calculation in
Chapter~\ref{ch:effective-calculus} proves
\[
 \int_a^b L(x)\,dx=[xL(x)-x]_a^b\qquad(0<a<b).
\]
Another is the exponential itself:
\[
 \int_a^b E(x)\,dx=E(b)-E(a).
\]
The left sides are rectangle computations and the right sides are finite
combinations of independently constructed values.

Finally define, by the factorial series, $E(z)$ for rational complex
inputs.  Coefficient separation gives
\[
 E(i\pi x)=C(x)+iS(x).
\]
For completeness, the even and odd factorial prefixes have derivatives
$C'=-\pi S$, $S'=\pi C$ and the correct initial values.  Equality with
the geometric pair follows from the following finite uniqueness argument:
the difference satisfies its integral equation, and iterating that
equation $n$ times bounds it by $B(\pi h)^n/n!$ on an interval of length
$h$.  The integral equations follow by telescoping the already proved
uniform increment bounds for both constructions.  This error tends to
zero.  Chapter~\ref{ch:differential-equations} gives the matrix form of
this same argument.
